\section{Minor Examination, 2020--21 (Semester 2)}

\subsection*{Q1. Process API}

\subsubsection*{(i) Child and parent after \texttt{fork()}}
\textbf{Question.} Explain the key difference between a child process generated
by the \texttt{fork()} function in C, and its parent. \hfill [2 points]

\textbf{Solution.}
Immediately after a successful \texttt{fork()}, the child is a nearly identical
copy of the parent: it has the same code and an initially identical memory
image, open-file descriptors, and register state.  The important differences
are:
\begin{itemize}
  \item they are different processes and therefore have different process IDs;
  \item \texttt{fork()} returns \(0\) in the child but returns the child's PID
        in the parent; and
  \item they have separate virtual address spaces.  Modern systems usually
        implement the initial copy using copy-on-write, but a later write by
        one process does not alter the other process's private memory.
\end{itemize}
Thus the return value is what normally lets the two processes follow different
branches after the fork.

\subsubsection*{(ii) Output of the fork/exec/wait program}
\textbf{Question.} Predict the output of following C program. Assume that
function \texttt{mywc()} takes a filename as parameter and writes the number of
words in it (integer) in standard output followed by filename. Justify your
answer. \hfill [3 points]
\begin{verbatim}
int main(int argc, char *argv[])
{
        printf("First: %d \n", (int) getpid());
        int rc = fork();
        if (rc == 0) {
                printf("Second: (pid:%d)\n", (int) getpid());
                char *myargs[3];
                myargs[0] = strdup("mywc");
                myargs[1] = strdup("abc.c");
                myargs[2] = NULL;
                execvp(myargs[0], myargs);
                printf("Second: %d"\n, (int) getpid());
        } else {
                wait(NULL);
                printf("First: %d\n Second:%d\n", rc, (int) getpid());
        }
return 0;
}
\end{verbatim}

\textbf{Solution.}
Strictly, the program as printed does not compile: in
\begin{verbatim}
printf("Second: %d"\n, (int) getpid());
\end{verbatim}
\verb|\n| lies outside the string literal. Thus the literal answer for the
printed program is a compiler diagnostic and no run-time output.

For the evidently intended line
\begin{verbatim}
printf("Second: %d\n", (int) getpid());
\end{verbatim}
let \(P\) be the parent's PID, \(C\) the child's PID, and \(W\) the integer word
count printed by \texttt{mywc} for \texttt{abc.c}. Assuming terminal
line-buffered output, if \texttt{execvp()} succeeds, the output is
\begin{verbatim}
First: P
Second: (pid:C)
W abc.c
First: C
 Second:P
\end{verbatim}
Here the last two numbers are not reversed accidentally: in the parent,
\texttt{rc} is the child's PID \(C\), whereas \texttt{getpid()} returns the
parent's PID \(P\).  The parent's \texttt{wait(NULL)} ensures that these last
two lines occur only after the child (and hence \texttt{mywc}) terminates.
The successful \texttt{execvp()} replaces the child's process image, so the
post-exec \texttt{printf()} is not executed.

If the corrected program's \texttt{execvp()} fails, there is no
\texttt{mywc} output and execution continues at the next statement in the
child:
\begin{verbatim}
First: P
Second: (pid:C)
Second: C
First: C
 Second:P
\end{verbatim}
If standard output is redirected and fully buffered, inherited user-space
buffers can change what is observed; the ordering above is the ordinary
terminal-output interpretation intended by the question.

\subsubsection*{(iii) Number of processes generated}
\textbf{Question.} How many new processes the following snippet will generate?
Why? \hfill [2 points]
\begin{verbatim}
for (i = 0; i < n; i++)
    fork();
\end{verbatim}

\textbf{Solution.}
Assuming all calls succeed, after each iteration every currently existing
process executes the next
\texttt{fork()}, so the population doubles.  Starting with one process, after
\(n\) iterations there are
\[
  2^n \quad\hbox{total processes}.
\]
Therefore the loop generates \(\fbox{\(2^n-1\)}\) new processes.  Equivalently,
the repeated forks form a binary process tree with \(2^n\) processes at the
end.

\subsection*{Q2. Multilevel Feedback Queue (MLFQ)}

\subsubsection*{(a) Approximating SJF}
\textbf{Question.} Explain how MLFQ approximates SJF without knowing how long
the jobs are going to run? \hfill [2 points]

\textbf{Solution.}
MLFQ infers a job's likely behavior from its observed CPU use.  New jobs begin
at high priority, normally with a short time slice.  A job that repeatedly
uses its complete allotment appears CPU-bound and is demoted to lower-priority
queues with longer slices.  A job that quickly yields or blocks for I/O
appears interactive or short and remains at high priority.  Thus short and
interactive jobs receive service before long CPU-bound jobs, approximating
SJF without requiring their lengths in advance.  Periodic priority boosts
prevent starvation and allow behavior changes to be recognized.

\subsubsection*{(b) Schedule with a 50 ms priority boost}
\textbf{Question.} Assume an MLFQ implementation with following parameters,
and that the context switch is instantaneous.
\begin{enumerate}
  \item[(a)] Number of queues: 4,
  \item[(b)] timer interrupt: 10ms,
  \item[(c)] Time slices for different queues (assume to be same as the time
        allotted for a process for that queue): 1st: 10ms, 2nd: 20ms,
        3rd: 30ms, 4th: 40ms,
  \item[(d)] Priority boost every 50ms.
\end{enumerate}
Describe with the help of a diagram (similar to ones done in class) which
processes get to run on the CPU at different times in the first 100 seconds in
the following scenario.
\begin{itemize}
  \item Process A enters at \(t=0\)ms and runs for 8ms before requesting for
        I/O, which takes 5ms to complete.
  \item Process B enters at \(t=30\)ms and needs to run for 100ms, and does not
        perform any I/O.
\end{itemize}
\hfill [4 points]

\textbf{Solution.}
The paper says ``first 100 seconds,'' but all arrivals, bursts, slices, and
boost periods are in milliseconds.  The dimensionally consistent intended
horizon is the first \(100\) ms.

The statement says when A's first I/O completes but does not say whether A
terminates, repeats the same CPU/I/O cycle, or has some other next CPU burst.
Consequently it does not determine a unique schedule.  One common classroom
interpretation is that A repeatedly runs for 8ms and then performs I/O for
5ms, and that voluntarily yielding before the end of a slice gives it a fresh
slice at the same priority.  With that convention, and with no
equal-priority arrival preemption, the \(0\)--\(100\)ms timeline is:
\[
\begin{array}{c|rrrrrrrr}
\text{interval (ms)}
&0\!-\!8&8\!-\!13&13\!-\!21&21\!-\!26&26\!-\!34&34\!-\!44
&44\!-\!52&52\!-\!62\\ \hline
\text{CPU}
&A&\text{idle}&A&\text{idle}&A&B(Q_1)&A&B(Q_1)
\end{array}
\]
\[
\begin{array}{c|rrrrrr}
\text{interval (ms)}
&62\!-\!70&70\!-\!75&75\!-\!83&83\!-\!88&88\!-\!96&96\!-\!100\\ \hline
\text{CPU}
&A&B(Q_2)&A&B(Q_2)&A&B(Q_2)
\end{array}
\]
At \(t=50\) the boost moves the waiting B from \(Q_2\) to \(Q_1\); A is
already running in \(Q_1\), so it continues until blocking at \(t=52\).  At
\(t=100\), B is boosted to \(Q_1\).  By then B has received \(34\)ms of its
\(100\)ms demand.

Under OSTEP's later anti-gaming Rule 4, CPU use is instead accumulated across
I/O until the level allotment is exhausted. Because the question says the
slice and allotment are equal, A then has only 2ms left in \(Q_1\) after its
first I/O and is demoted after using those 2ms. Assuming a boost puts the
already-waiting A before the interrupted B, the repeated-burst timeline is:
\[
\begin{array}{c|rrrrrrrr}
\text{interval (ms)}
&0\!-\!8&8\!-\!13&13\!-\!15&15\!-\!21&21\!-\!26&26\!-\!30
&30\!-\!40&40\!-\!44\\ \hline
\text{CPU}
&A(Q_1)&\text{idle}&A(Q_1)&A(Q_2)&\text{idle}&A(Q_2)
&B(Q_1)&A(Q_2)
\end{array}
\]
\[
\begin{array}{c|rrrrrrr}
\text{interval (ms)}
&44\!-\!50&50\!-\!58&58\!-\!68&68\!-\!70&70\!-\!90&90\!-\!96
&96\!-\!100\\ \hline
\text{CPU}
&B(Q_2)&A(Q_1)&B(Q_1)&A(Q_1)&B(Q_2)&A(Q_2)&B(Q_3)
\end{array}
\]
This is the OSTEP cumulative-allotment answer; B has received 50ms by
\(t=100\)ms.

If instead A has only the one stated CPU burst and terminates after its I/O,
the also-consistent timeline is
\[
\begin{array}{c|rrrrrrr}
\text{interval (ms)}
&0\!-\!8&8\!-\!30&30\!-\!40&40\!-\!50&50\!-\!60&60\!-\!80&80\!-\!100\\
\hline
\text{CPU}
&A&\text{idle}&B(Q_1)&B(Q_2)&B(Q_1)&B(Q_2)&B(Q_3).
\end{array}
\]
Here B has received \(70\)ms by \(t=100\)ms.  If ``100 seconds'' is read
literally, a complete schedule after \(100\)ms still depends on the missing
post-I/O behavior of A; B eventually completes and the remaining interval is
otherwise idle in the one-burst interpretation.

\subsubsection*{(c) Schedule with no priority boost}
\textbf{Question.} Modify the diagram in above part to show what would happen
if there was no priority boost? \hfill [2 points]

\textbf{Solution.}
Under the repeated-burst, fresh-slice-on-I/O interpretation above:
\[
\begin{array}{c|rrrrrrrr}
\text{interval (ms)}
&0\!-\!8&8\!-\!13&13\!-\!21&21\!-\!26&26\!-\!34&34\!-\!44
&44\!-\!52&52\!-\!57\\ \hline
\text{CPU}
&A&\text{idle}&A&\text{idle}&A&B(Q_1)&A&B(Q_2)
\end{array}
\]
\[
\begin{array}{c|rrrrrrr}
\text{interval (ms)}
&57\!-\!65&65\!-\!70&70\!-\!78&78\!-\!83&83\!-\!91&91\!-\!96
&96\!-\!100\\ \hline
\text{CPU}
&A&B(Q_2)&A&B(Q_2)&A&B(Q_2)&A
\end{array}
\]
B accumulates its \(20\)ms \(Q_2\) allotment over four 5ms runs and is demoted
to \(Q_3\) at \(t=96\).

With OSTEP cumulative allotments, the no-boost repeated-burst timeline is
\[
\begin{array}{c|rrrrrrrr}
\text{interval (ms)}
&0\!-\!8&8\!-\!13&13\!-\!15&15\!-\!21&21\!-\!26&26\!-\!30
&30\!-\!40&40\!-\!44\\ \hline
\text{CPU}
&A(Q_1)&\text{idle}&A(Q_1)&A(Q_2)&\text{idle}&A(Q_2)
&B(Q_1)&A(Q_2)
\end{array}
\]
\[
\begin{array}{c|rrr}
\text{interval (ms)}&44\!-\!64&64\!-\!70&70\!-\!100\\ \hline
\text{CPU}&B(Q_2)&A(Q_2)&B(Q_3).
\end{array}
\]
At \(t=70\), A has exhausted its cumulative \(Q_2\) allotment and joins
\(Q_3\) behind B; B runs through the horizon.

Under the one-burst/termination interpretation for A, the timeline is
\[
\begin{array}{c|rrrrrr}
\text{interval (ms)}
&0\!-\!8&8\!-\!30&30\!-\!40&40\!-\!60&60\!-\!90&90\!-\!100\\ \hline
\text{CPU}
&A&\text{idle}&B(Q_1)&B(Q_2)&B(Q_3)&B(Q_4).
\end{array}
\]
B then continues in \(Q_4\) and finishes at \(t=130\)ms.

\subsection*{Q3. Memory Virtualization}

\subsubsection*{(a) Multiple segments versus one base--bounds slot}
\textbf{Question.} Explain the advantages of segmentation using multiple
base-bounds pairs of registers over slotting approach with one pair of
base-bounds registers. \hfill [3 points]

\textbf{Solution.}
\begin{itemize}
  \item Code, heap, and stack can be placed independently in physical memory;
        physical memory need not contain one large contiguous slot for the
        whole virtual address space.
  \item The unused virtual gap between a growing heap and a growing stack need
        not consume physical memory, improving utilization.
  \item Heap and stack can grow independently (and in opposite directions),
        subject to relocation and available space.
  \item Each segment can have separate protection and sharing properties; for
        example, code may be read-only and shared while stack data is private.
\end{itemize}

\subsubsection*{(b) Address translation and instruction execution}
\textbf{Question.} Take an example of segmentation approach where the virtual
memory is 16KB (14 bits address space), and the base-bounds registers are as
following.
\begin{center}
\begin{tabular}{c|c|c|c}
\textbf{Segment} & \textbf{Base} & \textbf{Size (max 4K)}
& \textbf{Grows positive?}\\ \hline
Code\(_{00}\)  & 32K & 2K & 1\\
Heap\(_{01}\)  & 34k & 3K & 1\\
Stack\(_{11}\) & 44K & 2K & 0
\end{tabular}
\end{center}
Explain the address translations that happen while executing the following
instructions starting from the fetch at location 100 in virtual memory of the
process. Assume that the \texttt{ebx} register holds the the address 13K,
which in the process's virtual memory has an integer variable stored in it with
value 1234. Assume that the \texttt{ecx} register holds the the address 15K
which has number 2345 stored in it. The first two bits of the virtual address
determine the segment.
\begin{itemize}
  \item 100: Load the content of the address stored in \texttt{ebx} to
        \texttt{eax}.
  \item 104: Increment \texttt{eax} by 105.
  \item 107: Copy the content of \texttt{eax} into the address that is stored
        in \texttt{ebx}.
  \item 111: Load the content of the address stored in \texttt{ecx} to
        \texttt{eax}.
\end{itemize}
\hfill [4 points]

\textbf{Solution.}
The two selector bits leave a 12-bit raw offset \(r\), with
\(0\leq r<4\)KB.  In OSTEP's grow-negative convention the stack offset is
interpreted as a signed displacement
\[
 d=r-4\text{KB},
\]
and translation is \(\mathrm{PA}=\mathrm{base}+d\).  A stack access is valid
only when \(-\mathrm{size}\leq d<0\).  This subtraction of 4KB is essential;
using \(44\text{K}-r\) would reverse which of 13K and 15K is in the 2K stack.

The instruction addresses all select code (\(00\)); if reached, their fetch
translations are
\[
\begin{array}{c|c}
\text{instruction VA} & \text{instruction-fetch PA}\\ \hline
100 & 32\text{K}+100=32868\\
104 & 32\text{K}+104=32872\\
107 & 32\text{K}+107=32875\\
111 & 32\text{K}+111=32879
\end{array}
\]
where \(1\)K \(=1024\) bytes.  In the actual execution, however, only the
first fetch occurs:
\[
13\text{K}=11\,0100\,0000\,0000_2,\qquad
r=1\text{K},\qquad d=1\text{K}-4\text{K}=-3\text{K}.
\]
Because \(-3\text{K}<-2\text{K}\), the load through \texttt{ebx} violates the
stack bound and raises a \textbf{segmentation fault}.  Thus \texttt{eax} is not
loaded with 1234, and the instructions at 104, 107, and 111 are not executed.
The statement that a variable exists at virtual address 13K is inconsistent
with the printed 2K grow-negative stack segment.

For completeness, the separately considered \texttt{ecx} address would be
valid:
\[
15\text{K}=11\,1100\,0000\,0000_2,\qquad
r=3\text{K},\qquad d=-1\text{K},
\]
so it translates to \(44\text{K}-1\text{K}=43\text{K}=44032\) and would load
2345 if the instruction at 111 were reached.  Selector \(10\), absent from the
table, is invalid.

\subsection*{Q4. Free-Space Management}

\subsubsection*{(a) How \texttt{free()} obtains the allocation size}
\textbf{Question.} We saw that the \texttt{free()} call does not take any size
parameter. How does it know how much of memory it needs to free?
\hfill [2 points]

\textbf{Solution.}
An allocator stores hidden metadata in a header immediately before (or
otherwise associated with) each allocated payload.  The header records at
least the block size and may also contain a magic value, allocation status,
or links.  If \(p\) is the payload pointer returned by \texttt{malloc()},
\texttt{free(p)} moves backward by the known header size, reads the block
size, and returns that entire block to the free list, possibly coalescing it
with adjacent free blocks.

\subsubsection*{(b) First-fit, best-fit, and worst-fit}
\textbf{Question.} Suppose free-list for a process looks like following. Here
we care only about the size, so the addresses are not mentioned.
\[
  \mbox{head}\longrightarrow 15\mbox{KB}\longrightarrow10\mbox{KB}
  \longrightarrow40\mbox{KB}\longrightarrow8\mbox{KB}
  \longrightarrow\mathrm{NULL}.
\]
Suppose a request for 10KB of memory comes, followed by another request of
6KB. Explain how the free-list would look like after the first and the second
request for Best-Fit, Worst-Fit and First-Fit strategies. List one disadvantage
of each of these strategies. \hfill [6 points]

\textbf{Solution.}
Assume allocation is from the selected block and any remainder replaces that
block at the same free-list position; allocator-header and alignment overheads
are omitted.

\paragraph{First-fit.}
The first \(10\)-KB request uses the first sufficiently large block, namely
the \(15\)-KB block, leaving \(5\) KB:
\[
  \fbox{\(5\longrightarrow10\longrightarrow40\longrightarrow8\)}.
\]
For the \(6\)-KB request, \(5\) KB is too small, so the next \(10\)-KB block is
used and leaves \(4\) KB:
\[
  \fbox{\(5\longrightarrow4\longrightarrow40\longrightarrow8\)}.
\]
\emph{Disadvantage:} it can accumulate small fragments near the front of the
list, increasing fragmentation and future search time.

\paragraph{Best-fit.}
The first request uses the exact \(10\)-KB block:
\[
  \fbox{\(15\longrightarrow40\longrightarrow8\)}.
\]
For the second request, the smallest sufficient block is \(8\) KB, leaving
\(2\) KB:
\[
  \fbox{\(15\longrightarrow40\longrightarrow2\)}.
\]
\emph{Disadvantage:} it generally scans the whole list and tends to create
tiny unusable remainders (splinters).

\paragraph{Worst-fit.}
The first request uses the largest block, \(40\) KB, leaving \(30\) KB:
\[
  \fbox{\(15\longrightarrow10\longrightarrow30\longrightarrow8\)}.
\]
The second request again uses the largest block, now \(30\) KB, leaving
\(24\) KB:
\[
  \fbox{\(15\longrightarrow10\longrightarrow24\longrightarrow8\)}.
\]
\emph{Disadvantage:} it requires a full-list search (unless extra data
structures are maintained) and repeatedly consumes the largest blocks, which
can make later large requests impossible.
